\chapter[Sermon 28. Here Is Described Adoration and Praise Giving, or a Hymn Song unto Christ of the Beasts and Elders.]{Here Is Described Adoration and Praise Giving, or a Hymn Song unto Christ of the Beasts and Elders.}
\label{sermon:028}
\markright{Sermon 28. Here Is Described Adoration and Praise Giving, or a ...}

And when he had taken the book, the four living creatures and the twenty-four elders fell down before the Lamb, having harps and golden vials full of fragrant incense (which are the prayers of the saints), and they sang a new song, saying: “You are worthy to take the book, and to open the seals thereof: for you were slain, and you have redeemed us by your blood, out of every tribe and tongue, and people and nation, and you have made us to our God, kings and priests, and we shall reign on the earth.”

We have heard that the Lamb has received the book from the hand of him who sits on the throne, so that he might open it and break the seals of it; that is, we understand that Christ is the only and eternal Savior and Lord, to whom all power is given in heaven and earth; that he alone and forever saves, that he reveals to us the mysteries and judgments of God, and that he finally governs and disposes all things in the world. It follows moreover how all the creatures of God behaved themselves toward this Son of God, the monarch and governor of all things. This is set forth with a marvelous figurative and plentiful speech in the vision of the four beasts, twenty-four elders, and so on. Certainly, that we might from their gestures, words, and works understand what is fitting for us to do in the judgments of God.

\index{Jerome, Saint}For this example is truly manifold, and even consists of six parts, such as you will scarcely find presented in any other matter. And this matter is of very great force. First, we have heard in the fourth chapter that the four beasts cried out before the throne of him who sat: “Holy, holy, holy, Lord God Almighty.” Second, we understand that the twenty-four elders fell down, worshipped, and cast away their crowns, and sang a hymn. Now follows the third degree of this example. For as first the beasts and elders did these things separately, so now jointly with one accord, the beasts and elders fall down together before the Lamb. Let us therefore fall down also in all the judgments and works of God, before the Lamb, governor of all, and let us worship. For although it is not here added, “and they worshipped,” they are to be understood to have fallen down for the purpose of worship, for to fall down is to worship. This is also evident from what follows. For they offer prayers to the Lamb, that is to say, sing a hymn, which is a part of godly worship. Moreover, it follows immediately that every creature sang a hymn to him who sits on the throne and to the Lamb, and so forth. And truly, two things especially and diligently Jerome treats in this example. First, he depicts the behavior of the beasts and elders. After that, he adds the hymn, praise-giving, or song. And so much as pertains to their behavior: before all things, they fall down before the Lamb, as I said even now.

And this passage is sufficiently powerful to demonstrate the deity of our Savior Christ. These points should be compared with what is written on the same topic in chapter four. The twenty-four elders fell down before him who sits on the throne, and worship him who lives forever and ever; and now it is said that the very same elders have fallen down before the Lamb. Consequently, he who lives forever and the Lamb are worshipped with equal glory, cult, and honor, and that the Son is coequal with the Father, to be worshipped eternally. Thus, the abominable and detestable error of Arius and Servetus is now openly perceived, refuted not only by beasts but also by the entire congregation of saints in heaven. Idle men reason subtly, perverting and twisting God’s word according to their habitually bold inclinations, at their pleasure. We will rather follow the examples of all saints and creatures in the world, and will worship the Lamb with him who sits on the throne, blessed forevermore.

\index{John, Saint}\index{Arethas of Caesarea}Again, there are objections to us: the Elders lying prostrate on the pavement, holding in their hands harps and vials. A harp in the Psalms and holy history is an instrument of music, consecrated to divine praises. Concerning the vial, writers of vessels discuss much about its shape and fashion; I understand it to be simply a cup or a bowl, such as we read there were many in the tabernacle and temple, appointed both for drink offerings and also for sweet odors and incense. Nevertheless, these things in the holy heavenly dwellers are not to be taken corporally, but spiritually, after a figuration. For what the Spirit of God understood, the revealer of secrets, Saint John himself adds, which are the prayers of Saints. Therefore, it is signified that Saints offer prayers to God, which are much more acceptable to him than the sweet melody of musical instruments is to man, or the pleasant savor of sweet gums or of incense. Arethas, the expositor, says that because they have harps, it shows a concord and agreement in giving God thanks. And from this we learn again what we should do in the contemplation and understanding of the judgments and works of God. The Lord is to be praised and blessed because he is good, and his mercy endures forever. But if thanks must be given to God, if his works and judgments are to be praised, why do certain men expostulate with God, blaming or bringing into suspicion his judgments? Let us learn moreover that organs and those corporal incenses do no longer become the church of God.

\index{Irenaeus}\index{Tertullian}\index{Papacy}Saint Irenaeus, in book 4 against heresies, chapters 33 and 34, demonstrates that the prayers and thanksgiving of the saints are the same offering that Malachi prophesied would be offered throughout the world. Shortly afterward, Tertullian followed the same interpretation against the Jews and in book 4 against Marcion; other doctors of the Church have followed them. But those pleasant, sophistical triflers, I mean the papist divines, as if triumph in these things, yet they lead a shadowing and a most vain triumph. For they apply these things to their sacrifice, wherein they pretend under the guise of bread and wine to offer up to God the Father the body and blood of Christ, a propitiatory sacrifice for the living and the dead. But Irenaeus and Tertullian speak not of such a sacrifice, but of the oblation of prayers, which the mass-priest offers not up alone, but the whole congregation of Christ sanctified in his blood, giving thanks in the Lord’s Supper to God the Father for their free redemption. These holy fathers never knew the sale Masses of these Canaanites.

\index{Augustine, Saint}Regarding this same passage from the writings of John, the Roman Catholics attempt to use it to prove and establish the practice of praying to saints in heaven. “Behold,” they say, “the saints are said to pray openly in heaven.” But they argue that they do not pray for themselves, and therefore, as intercessors and patrons, they pray for their clients and worshippers on earth. I reply that the saints indeed pray in heaven, but you, by adding your own explanations of the kind and manner of that prayer, distort it, fabricate, and maliciously and falsely misrepresent it. John explains himself here, and does not require your additions. For he adds, “and they sing a new song.” Indeed, he recites the entire form of this song, lest anyone corrupt what he had said about prayers. And that form contains praise and blessing or thanksgiving, and not intercession or invocation. It is certain, even according to the Apostle’s doctrine, in 1 Timothy 2 and Philippians 4, that there are two chief kinds of prayer: invocation and praise or thanksgiving. But the passage itself plainly shows that John speaks here of the latter, not the former. And just as this passage explains certain types, shadows, or mysteries of God’s law, we can more effectively refute the intercession of saints in heaven for their worshippers. For in the law, only one golden altar of incense was permitted. And that altar represented a figure of Christ. For one Christ is the mediator and intercessor between God and man. It was not lawful for the people of God to burn incense, but only upon this altar. It was not lawful for any man to prepare or make for himself an odor of those kinds of gums from which the divine incense consisted, and to offer it to God, as appears in Exodus 30. Why, then, do these people not understand that prayers belong to God alone, and that the saints in heaven would not offer such incense? David, in Psalm 141, says, “Let my prayer be directed as incense in your sight, the lifting up of my hands an evening sacrifice.” The Devil desires to have such incense made to him, as appears in Matthew 4 and in Augustine’s \textit{The City of God}. But our heavenly saints are not devils. Why do they not understand that this altar of incense now stands in heaven on the right hand of the Father, and there makes intercession for us? And that for his sake the Father is reconciled to us, and we are accepted of God, and that through him alone we must offer up our prayers to God, which are otherwise abominable? Why do they not see the heavenly saints at this present attributing all things to the only Lamb alone, and claiming nothing for themselves? Finally, that they make no mention of their worshippers, but plainly testify that the only Lamb was and is worthy, who should take the book, etc.?

And the praise or thanksgiving of the heavenly saints he has called a new song, which in the Scriptures is no new thing. For the saints say that they will sing on Earth to God a new song, Psalm 33, 96, 98, 149, Isaiah 42. And new songs are called these new ballads or verses in meter, which are made of some new benefit or noble act done. And because the mind of man is greatly delighted with new benefits, they sing a new song, which with a joyful mind praises God and gives him thanks with their inward affections. Finally, they sing a new song, which with purified minds and renewed by the Spirit does laud God; which thing was chiefly given to those heavenly saints. From this we learn again how our minds ought to be affected and furnished in the prayers and praises of God. This same, says Arethas, I call a new song, by whose benefit we, who being enlightened in all parts of the earth, departing from the antiquity of the law written, and walking in the newness of life, are taught by the Holy Spirit to sing a giving of thanks.

To these things now is added the Hymn of Saints, so that we might also have a form whereby to praise God. And in the Hymn they sing that all things are subject to Christ, and all things ordered by his rule, that he humbled himself to death, and was therefore exalted above all things. Now also the virtues or effects and wonderful benefits of his death are commended to us, that, esteeming the governor of the benefits done to us, we may believe also that his rule will be wholesome for us, and therefore may submit ourselves to him willingly in faith and patience. Which indeed is the chiefest end of those things which here are treated with such diligence.

\index{Resurrection}First, they commend the majesty and dignity of Christ, stating that he alone is found throughout the entire world, possessing dominion over all, the sole savior of the world, the revealer of divine mysteries, and the governor of all. For this is to take and to open the book, which we have often repeated. Secondly, they add the reason why this glory should belong only to the Lamb or Son of God: because, they say, you were slain. And they understand by the lesser the greater, namely, his entire incarnation and the entire mystery of our redemption, death, resurrection, and ascension into heaven, and the remainder. Therefore, he is the true and only mediator of God and men; he is the only savior, as he alone was incarnate and crucified for us; he is the only governor, who, through his humility, deserved to be exalted, Philippians 2. And he is a most suitable governor of all things, as from whom all men may, as their most faithful savior and even their brother, hope well, whatever events befall them through his governance, etc.

\index{Hebrews, Epistle to the}Meanwhile, they commend highly the power or effect of Christ’s death. For if this is rightly understood, we are more ready to submit ourselves to that governor whom we know to be our savior, who loves us dearly and desires all to be saved. And the chiefest effect of Christ’s death is redemption; redemption includes deliverance from captivity. We were prisoners and servants of sin, of death, and utter slaves of the devil and hell. And the Son of God came and took flesh, and shed his blood (for so also is the manner of redeeming us expressed by the elders), and he has washed us from our sins, and being purified, he has ransomed us from the power of death, hell, sin, and Satan, so that now we belong to God. Therefore, they say expressly, “You have redeemed us to God.” We therefore belong to God; the devil has no more power over us; we are the freemen of Christ, delivered through his blood. 1 Peter, 1; Hebrews, 9. And since we now belong to God, namely, justified freely by his grace, through the blood of Christ, as the apostle also says in the 3rd chapter to the Romans, we ought to serve God truly in the newness of spirit, not the flesh and the devil, in the oldness of the letter and of our flesh. Which the same apostle discusses more at length in the 6th chapter to the Romans.

\index{Daniel (Book of)}Those whom he has redeemed declare it openly, men truly of every tribe, and so forth. In this recounting, he imitates Daniel in chapter 7, signifying universality. For the Lord died for all; but that all are not made partakers of this redemption is through their own fault. For the Lord excludes no one, but only him who, through his own unbelief, excludes himself.

Following redemption is another effect of Christ's death, for it makes men justified to God as kings and priests. For those who are justified, work righteousness. I have explained this passage concerning the priesthood and kingdom of Christians in the first chapter, where you may find it.

\index{Turks}The saints also affirm that they will reign on earth, namely through the power of Christ; not physically, as the Millenarians imagine, and the Turks following the same, imagining physical things in this world and joys in a terrestrial paradise. For the entire Scripture promises better things. Nor must the godly be so devoted to physical things that they should hope for nothing beyond physical matters. The saints speak here of the final judgment, in which it will be shown to the entire world and to all who dwell on earth that the saints, who at one time seemed to the world to have been wicked, ungodly, peacebreakers, heretics, and parricides, and for that reason were slain, are righteous, holy kings and priests of God. Thus I say they will reign on earth. This matter is explained more fully in chapters 3 and 5 of the Book of Wisdom.

Let the faithful consider these things, I say, when they are oppressed by the wicked for the truth and righteousness, through the permission of Christ, the governor of all, in this world; nevertheless, let them glorify the Lord God and praise him without ceasing. To him be glory forever.
